% § 8 — Results C: When a Query Fires and the Benchmark Does Not
\section{Results C: When a Query Fires and the Benchmark Does Not}
\label{sec:audit}
Every flag returns the template fields that satisfied the query and the sentences they were extracted from, so each error can be traced to one of three sources: a mis-filled field, a query broader than its grey-list item, or a missing reference label for a clause that a literal reading of the Annex captures. In this study the attribution is automatic and partial. \textbf{Misses.} No missed sentence lacks a template altogether: every one lies in a clause for which templates were extracted. For unilateral change, most misses carry no template with the item's action. For limitation of liability and termination, about half carry the right action but a field that blocks the query, most often an exclusion whose object names neither personal injury nor gross negligence (42 sentences for \gi{a} and \gi{b}), or a termination read as requiring cause (15 for \gi{g}). \textbf{False positives.} Item \gi{q} accounts for 86 of the 162 item-level false positives, all in dispute-resolution clauses that \corpus{} leaves unlabelled at the flagged sentence. Across scored items, 66 item-level false positives lie in a clause whose other sentences carry the mapped label, which points to the sentence-level granularity of the reference rather than to the query. Finally, the items without a \corpus{} counterpart fire 77 times: 39 flags for \gi{m}, none of them labelled, 17 for \gi{p} and 21 for \gi{d}. Deciding whether such flags are grounded requires legal judgement; it is left to the expert study described in Section~\ref{sec:discussion}.
